%% kapa: exam=1 round=27 year=2016 period=1 form=A subject=civil_law subject_ko=민법 seq=20 number=20
%% source: Q-Net 감정평가사 기출문제 (cst00309) artlSeq=5169402 file=extracted/1008/5169402/제27회 감정평가사 1차 문제/제27회 감정평가사 1교시문제(A형).pdf
%% conversion: kordoc 4.24.51 -> tools/gen1.py (auto); review=proofread-sonnet; flags=none
\begin{question}{20}
소멸시효에 관한 설명으로 옳지 않은 것은? (다툼이 있으면 판례에 따름)
\begin{choices}
  \item 소유권에 기한 물권적 청구권은 소멸시효의 대상이 되지 않는다.
  \item 공유관계가 존속하는 한 공유물분할청구권만이 독립하여 시효로 소멸될 수 없다.
  \item 소멸시효를 주장하는 자가 제기한 소에 권리자가 응소하여 적극적으로 권리를 주장하고 그것이 받아들여진 경우, 응소한 때에 소멸시효가 중단된다.
  \item 근저당권설정등기청구의 소제기에는 그 피담보채권이 될 채권에 대한 소멸시효 중단효력은 없다.
  \item 소멸시효의 중단사유로서의 승인은 소멸시효의 진행이 개시된 이후에만 가능하다.
\end{choices}
\end{question}
\answer{4}
